\section{A joint physical and nonconvex certificate theorem}
\label{sec:v16-joint}
The finite testing and microscopic approximation results now give a
single lower--upper certificate for the original resource problem.
The stateless transfer inequality and physical-to-rational bridge
are those of Theorem~\ref{thm:v15-main}; the polynomial lower module
is a separate strengthening, not a new physical transfer principle.
Unlike a constant mixture of tests, the polynomial lower certificate
can separate a target from a private behavior image even when the
target belongs to its convex hull.

\begin{theorem}[Constructive same-budget physical testing]
\label{thm:v16-main}
Fix a finite horizon and one of the private, hidden-selector, or
visible-selector signatures. Let $E,F$ be marked physical acquisitions
with finite compatible controls and marks. Suppose that rational finite
presentations $E_j,F_j$ have executable two-sided stateless marked
feedback errors $a_j,c_j\to0$, with computable upper bounds. For the algorithmic assertion require
the uniform enumeration and rational error enclosures of
Proposition~\ref{prop:v17-effective}, not merely individual
computability of each presentation. Let
$P_{p,L}$ be the polynomial \eqref{eq:v16-positivepoly} of the finite
intrinsic deficiency problem for $E_j,F_j$, and let a nonnegative
rational Bernstein array certify it. Let an actual rational simulator
table at the same signature have finite error $U$. Then
\begin{equation}\label{eq:v16-main}
 \max\{0,L-a_j-c_j\}\le\delta_{\mathfrak b}(E,F)
                \le\min\{1,U+a_j+c_j\}.
\end{equation}
There is a terminating dovetailed search for such an interval of any
prescribed positive rational width. Every lower certificate is checked
by a finite polynomial identity and sign comparisons; every upper
certificate is an executable table at the original budget. The visible
case retains the selector in both compared marked laws.

The hypotheses hold for the effective fixed-particle collision class
of Theorem~\ref{thm:v15-constructive-core}, and in particular for the
nonconserved collective acquisitions of Theorem~\ref{thm:v16-taylor}
with equilibrium preparation and a finite computable marked channel
with the piecewise effective integration property proved there.
For the sign mark of Theorem~\ref{thm:v16-exactphysical}, that integration
property is explicit.
\end{theorem}
\begin{proof}
The finite polynomial identity gives
$L\le\delta_{\mathfrak b}(E_j,F_j)$ by
Theorem~\ref{thm:v16-dual}; direct evaluation of the rational table
gives the reverse bound by $U$. Compose it with the two stateless
physical presentation channels. Their identity state costs preserve
all retained widths and selector budgets. The marked triangle and
presentation-stability inequalities give
\[
 |\delta_{\mathfrak b}(E,F)-\delta_{\mathfrak b}(E_j,F_j)|
                       \le a_j+c_j.
\]
This proves \eqref{eq:v16-main}; no attainment assertion for the original
nonfinite minimum is required. The visible comparison is joint with
the same selector, not its hidden marginal.

Given a target width $\eta>0$, enumerate presentations until their
certified sum of errors is less than $\eta/4$. For that presentation,
Corollary~\ref{cor:v16-algorithm} produces lower and upper endpoints
with gap less than $\eta/2$. Retain maxima of all verified lower
endpoints and minima of all verified upper endpoints. The resulting
nested physical intervals have width less than $\eta$ and converge.
Zero lower bounds are always available, so this argument also covers
zero deficiency and does not ask a strict-positivity certificate to
prove equality at a boundary.

The cited constructive collision theorem provides effective graph
approximation and marked integration; the finite-report and joint-prefix
rationalization theorems then give the presentations and their error
bounds. In the collective class, \eqref{eq:v16-taylor-error} replaces
the unquantified graph approximation by an explicit rate. Its compact
cylinder integration and Gaussian tail bounds give every finite
marked-prefix probability. The retained finite-presentation theorem
rationalizes these joint prefixes without dividing by a small
conditional probability. Thus all links in this specialization are
supplied either by an explicitly cited retained theorem or by a proof
above, not by assuming a rational physical transition oracle.
\end{proof}

\subsection{Exponential acquisition values}
The distinction between small total variation and accuracy at an
exponential scale is important for the nonlinear-observable direction
of the program. The following estimate gives its precise price. It
is a theorem about bounded marked acquisition functionals, not a
construction of a kinetic large-deviation action.

\begin{theorem}[Sharp exponential transfer]
\label{thm:v16-exponential}
Let $P,Q$ be two laws of the same marked transcript with
$\|P-Q\|_{\rm TV}\le\varepsilon$. If $h$ is measurable and
$A\le h\le B$, then for every $\beta>0$,
\begin{equation}\label{eq:v16-logtransfer}
 \left|\frac1\beta\log\mathbb E_Pe^{\beta h}
       -\frac1\beta\log\mathbb E_Qe^{\beta h}\right|
 \le\frac1\beta\log\{1+\varepsilon(e^{\beta(B-A)}-1)\}.
\end{equation}
The bound is sharp. If a pair of stateless marked presentation channels
has error $\varepsilon$ uniformly over a declared resource signature,
the payoff is evaluated on the same compared marked output interface.
Under this convention the same bound holds between optimal infimum values, and between
the optimal supremum values, over that signature. The policies are
transported through the channels at unchanged budgets.

For the effective physical class of Theorem~\ref{thm:v16-main}, let
$\beta_n>0$ and an upper bound $M_n\ge B_n-A_n$ be computable.
One can construct finite rational presentations with certified errors
$\varepsilon_n\le e^{-\beta_nM_n}/n$. Their optimized normalized
logarithmic acquisition values then differ from the physical ones
by at most $\log(1+1/n)/\beta_n$.
\end{theorem}
\begin{proof}
For any bounded $g$, integration of its level sets gives
$|\mathbb E_Pg-\mathbb E_Qg|\le\varepsilon(\sup g-\inf g)$.
Take $g=e^{\beta(h-A)}\in[1,e^{\beta(B-A)}]$. Both expectations
are at least one. Their larger-to-smaller ratio is at most
$1+\varepsilon(e^{\beta(B-A)}-1)$, which proves the estimate.
Equality holds when $Q$ is concentrated where $h=A$ and $P$ moves
mass $\varepsilon$ to a point where $h=B$.

For any admissible physical policy transport it to the presentation
and apply the uniform inequality; this gives one inequality between
optimal values. Transport in the other direction for the reverse
inequality. Approximate optimizing policies suffice, so attainment is
unnecessary. The same reasoning works for infima and suprema, but does
not replace a common policy by separate phasewise optimizers.

Finally choose a strictly positive rational tolerance below
$e^{-\beta_nM_n}/n$, using computable exponential enclosures. The
terminating presentation search of Theorem~\ref{thm:v16-main} attains
that tolerance. Substitution gives the final claim. This is an
accuracy-versus-scale construction; it asserts neither an efficient
complexity rate nor a large-deviation limit of the microscopic model.
\end{proof}

\begin{remark}[The downstream implication]
Theorems~\ref{thm:v16-collective}--\ref{thm:v16-exponential} provide a
verified fixed-particle link from a nonconserved collision-domain
observable to finite-resource marked and exponential acquisition
values, with rigorously computable errors. They do not assert
particle-number-uniform bounds, a Boltzmann--Grad limit, a nonlinear
BBGKY graph corrector on action sublevels, or a range theorem for the
historical kinetic Hamiltonian. Those are distinct analytic assertions
in the retained program. In particular \eqref{eq:v16-logtransfer}
explains why plain unscaled total variation alone does not supply
exponential-scale control. All eleven historical components remain in
the dependency graph; neither component order nor the existence of an
adapter is used as a proof of a missing kinetic theorem.
\end{remark}
